\documentclass[11pt,a4paper]{article}
\usepackage[utf8]{inputenc}
\usepackage[spanish]{babel}
\usepackage{amsmath, amssymb}
\usepackage{geometry}
\geometry{top=2.5cm, bottom=2.5cm, left=2.5cm, right=2.5cm}

\begin{document}

\section*{Piense, dialogue y comparta}

\begin{itemize}
    \item[\textbf{1.}] A un estudiante se le suministra un gran número de papel de copia, regla, compás, tijeras y una balanza sensible. Corta varias formas en diferentes tamaños, calcula sus áreas, mide sus masas y traza la gráfica de la figura TP1.1. (a) Considere el cuarto punto experimental de la parte superior. ¿Qué tan lejos está verticalmente de la recta de mejor ajuste? Exprese su respuesta como diferencia en la coordenada del eje vertical. (b) Exprese su respuesta como un porcentaje. (c) Calcule la pendiente de la recta. (d) Indique lo que la gráfica muestra, refiriéndose a la forma de la gráfica y los resultados de los incisos (b) y (c). (e) Describa si este resultado debe ser previsto teóricamente. (f) Describa el significado físico de la pendiente.
    \vspace*{12cm}
    
    \item[\textbf{2.}] \textbf{ACTIVIDAD} Cada persona en el grupo mide la estatura de otra persona utilizando una regleta graduada con distancias métricas en un lado y las distancias usuales de Estados Unidos, como pulgadas, en el otro lado. Registre la estatura al centímetro más cercano y a la media pulgada más cercana. Para cada persona, divida su estatura en centímetros entre la estatura en pulgadas. Compare los resultados de esta división para todos en su grupo. ¿Qué puede decir acerca de los resultados?
    \vspace*{6cm}
    
    \item[\textbf{3.}] \textbf{ACTIVIDAD} Reúna un número de centavos de Estados Unidos, ya sea de su profesor o de los miembros de su grupo. Divida los centavos en dos muestras: (1) los de fechas de 1981 o anteriores, y (2) los de fechas de 1983 y posteriores (excluya los centavos de 1982 de su muestra). Encuentre la masa total de todos los centavos en cada muestra. Luego divida cada una de estas masas totales entre el número de centavos en su muestra correspondiente, para encontrar la masa media de los centavos en cada muestra. Explique por qué los resultados son diferentes para las dos muestras.
    \vspace*{6cm}
    
    \item[\textbf{4.}] \textbf{ACTIVIDAD} Discuta en su grupo el proceso por el cual usted puede obtener la mejor medida del espesor de una sola hoja de papel en los capítulos 1 a 5 de este libro. Realice la medición y exprésela con un número adecuado de cifras significativas e incertidumbre. De esa medida, prediga el espesor total de las páginas de este libro (capítulos 1-10). Después haga su predicción, mida el grosor del volumen 1. ¿Está su medida dentro del rango de predicción y su incertidumbre asociada?
    \vspace*{6cm}
\end{itemize}

\newpage

\section*{Problemas}

\subsection*{SECCIÓN 1.1 Estándares de longitud, masa y tiempo}

\begin{itemize}
    \item[\textbf{1.}] (a) Use la información que aparece al final de este libro para calcular la densidad promedio de la Tierra. (b) ¿Dónde encaja el valor entre los que se mencionan en la tabla 10.1 en el capítulo 10? Busque la densidad de una roca superficial típica, como el granito, en otra fuente y compare la densidad de la Tierra con ella.
    \vspace*{7cm}
    
    \item[\textbf{2.}] Un protón, que es el núcleo de un átomo de hidrógeno, se representa como una esfera con un diámetro de $2.4 \text{ fm}$ y una masa de $1.67 \times 10^{-27} \text{ kg}$. (a) Determine la densidad del protón. (b) Establezca cómo se compara su respuesta del inciso (a) con la densidad del osmio, que está dada en la tabla 10.1 en el capítulo 10.
    \vspace*{7cm}
    
    \item[\textbf{3.}] De cierta roca uniforme se cortan dos esferas. Una tiene $4.50 \text{ cm}$ de radio. La masa de la segunda esfera es cinco veces mayor. Encuentre el radio de la segunda esfera.
    \vspace*{5cm}
    
    \item[\textbf{4.}] ¿Qué masa se requiere de un material con densidad $\rho$ para hacer un cascarón esférico hueco que tiene radio interior $r_1$ y radio exterior $r_2$?
    \vspace*{4cm}
    
    \item[\textbf{5.}] Usted ha sido contratado por el abogado defensor como un testigo experto en una demanda. El demandante es alguien que acaba de ser un pasajero en el primer vuelo de turista espacial orbital. Basado en un folleto de viajes ofrecido por la compañía de viajes espaciales, el demandante esperaba ver la gran muralla de China desde su altura orbital de $200 \text{ km}$ por encima de la superficie de la Tierra. No pudo hacerlo y ahora está exigiendo que se reembolse la tarifa y se le dé una compensación financiera adicional para cubrir su gran decepción. Construya la base para un argumento para la defensa que muestra que su expectativa de ver la gran muralla desde la órbita no era razonable. La pared es de $7 \text{ m}$ en su punto más ancho y la agudeza visual normal del ojo humano es $3 \times 10^{-4} \text{ rad}$. (La agudeza visual es el ángulo más pequeño subtendido que un objeto puede hacer en el ojo y aún reconocerse; el ángulo subtendido en radianes es el cociente del ancho de un objeto entre la distancia del objeto desde sus ojos.)
    \vspace*{8cm}
\end{itemize}

\newpage

\subsection*{SECCIÓN 1.2 Modelado y representaciones alternativas}

\begin{itemize}
    \item[\textbf{6.}] Una topógrafa mide la distancia a través de un río recto con el siguiente método (figura P1.6). Comenzando directamente a través de un árbol en la orilla opuesta, camina $d = 100 \text{ m}$ a lo largo de la ribera para establecer una línea de base. Luego mira al otro lado del árbol. El ángulo desde su línea de base hasta el árbol es $\theta = 35.0^\circ$. ¿Qué tan ancho es el río?
    \vspace*{6cm}
    
    \item[\textbf{7.}] Un sólido cristalino consta de átomos configurados en una estructura reticular repetitiva. Considere un cristal como el que se muestra en la figura P1.7a. Los átomos residen en las esquinas de los cubos de lado $L = 0.200 \text{ nm}$. Una pieza de evidencia para el ordenamiento regular de átomos proviene de las superficies planas a lo largo de las cuales se separa o fractura un cristal cuando se rompe. Suponga que este cristal se fractura a lo largo de una cara diagonal, como se muestra en la figura P1.7b. Calcule el espaciamiento $d$ entre dos planos atómicos adyacentes que se separan cuando el cristal se fractura.
    \vspace*{7cm}
\end{itemize}

\subsection*{SECCIÓN 1.3 Análisis dimensional}

\begin{itemize}
    \item[\textbf{8.}] La posición de una partícula que se mueve con aceleración uniforme es una función del tiempo y de la aceleración. Suponga que escribimos esta posición $x = k a^m t^n$, donde $k$ es una constante adimensional. Demuestre usando análisis dimensional que esta expresión se satisface si $m = 1$ y $n = 2$. ¿Puede este análisis dar el valor de $k$?
    \vspace*{6cm}
    
    \item[\textbf{9.}] ¿Cuáles de las siguientes ecuaciones son dimensionalmente correctas? (a) $v_f = v_i + ax$ (b) $y = (2 \text{ m}) \cos(kx)$, donde $k = 2 \text{ m}^{-1}$.
    \vspace*{6cm}
    
    \item[\textbf{10.}] (a) Suponga que la ecuación $x = At^3 + Bt$ describe el movimiento de un objeto en particular, siendo $x$ la dimensión de la longitud y $t$ la dimensión de tiempo. Determine las dimensiones de las constantes $A$ y $B$. (b) Determine las dimensiones de la derivada $dx/dt = 3At^2 + B$.
    \vspace*{6cm}
\end{itemize}

\newpage

\subsection*{SECCIÓN 1.4 Conversión de unidades}

\begin{itemize}
    \item[\textbf{11.}] Una pieza sólida de plomo tiene una masa de $23.94 \text{ g}$ y un volumen de $2.10 \text{ cm}^3$. A partir de estos datos, calcule la densidad del plomo en unidades SI ($\text{kg/m}^3$).
    \vspace*{5cm}
    
    \item[\textbf{12.}] ¿Por qué no es posible la siguiente situación? El dormitorio de un estudiante mide $3.8 \text{ m}$ por $3.6 \text{ m}$, y su techo tiene $2.5 \text{ m}$ de altura. Después de que el estudiante complete su curso de física mostrará su dedicación empapelando por completo las paredes de la habitación con las páginas de este libro de texto. Incluso cubrirá la puerta y la ventana.
    \vspace*{7cm}
    
    \item[\textbf{13.}] Un metro cúbico ($1.00 \text{ m}^3$) de aluminio tiene una masa de $2.70 \times 10^3 \text{ kg}$, y el mismo volumen de hierro tiene una masa de $7.86 \times 10^3 \text{ kg}$. Encuentre el radio de una esfera de aluminio sólida que equilibraría una esfera de hierro sólida de $2.00 \text{ cm}$ de radio sobre una balanza de brazos iguales.
    \vspace*{6cm}
    
    \item[\textbf{14.}] Sea $\rho_{Al}$ la densidad del aluminio y $\rho_{Fe}$ la del hierro. Encuentre el radio de una esfera de aluminio sólida que equilibra una esfera de hierro sólida de radio $r_{Fe}$ sobre una balanza de brazos iguales.
    \vspace*{4cm}
    
    \item[\textbf{15.}] Un galón de pintura (volumen $= 3.78 \times 10^{-3} \text{ m}^3$) cubre un área de $25.0 \text{ m}^2$. ¿Cuál es el grosor de la pintura fresca sobre la pared?
    \vspace*{5cm}
    
    \item[\textbf{16.}] Un auditorio mide $40.0 \text{ m} \times 20.0 \text{ m} \times 12.0 \text{ m}$. La densidad del aire es $1.20 \text{ kg/m}^3$. ¿Cuáles son (a) el volumen de la habitación en pies cúbicos y (b) el peso en libras del aire en la habitación?
    \vspace*{7cm}
\end{itemize}

\newpage

\subsection*{SECCIÓN 1.5 Estimaciones y cálculos de orden de magnitud}

\begin{itemize}
    \item[\textbf{17.}] (a) Calcule el orden de magnitud de la masa de una bañera medio llena de agua. (b) Calcule el orden de magnitud de la masa de una bañera medio llena de monedas de cobre.
    \vspace*{7cm}
    
    \item[\textbf{18.}] En un orden de magnitud, ¿cuántos afinadores de piano residen en la ciudad de Nueva York? El físico Enrico Fermi fue famoso por plantear preguntas como ésta en los exámenes orales para calificar a los candidatos a doctorado.
    \vspace*{8cm}
    
    \item[\textbf{19.}] Su compañero de cuarto está jugando un videojuego de la última película de \textit{Star Wars} mientras usted estudia física. Distraído por el ruido, usted va a ver lo que está en la pantalla. El juego implica tratar de volar una nave espacial a través de un campo lleno de asteroides en el cinturón de asteroides alrededor del Sol. Usted le dice: ``¿Sabes que el juego que estás jugando es muy poco realista? ¡El cinturón de asteroides no está tan lleno de gente y no tienes que maniobrar de esa manera!''. Distraído por su declaración, él accidentalmente permite que su nave espacial le pegue a un asteroide, justo para no lograr una puntuación alta. Se vuelve hacia usted con disgusto y dice: ``Sí, demuéstralo''. Usted dice: ``Bien, he aprendido recientemente que la mayor concentración de asteroides está en una región con forma de rosquilla entre los huecos de Kirkwood en los radios de $2.06 \text{ UA}$ y $3.27 \text{ UA}$ del Sol. Se estima que hay $10^9$ asteroides con un radio de $100 \text{ m}$ o más grande, como los de tu videojuego, en esta región\dots''. Termine su argumento con un cálculo para demostrar que el número de asteroides en el espacio cerca de una nave espacial es pequeña. (Una unidad astronómica, UA, es la distancia media de la Tierra al Sol: $1 \text{ UA} = 1.496 \times 10^{11} \text{ m}$.)
    \vspace*{10cm}
\end{itemize}

\newpage

\subsection*{SECCIÓN 1.6 Cifras significativas}

\begin{itemize}
    \item[\textbf{20.}] ¿Cuántas cifras significativas hay en los siguientes números: (a) $78.9 \pm 0.2$ (b) $3.788 \times 10^9$ (c) $2.46 \times 10^{-6}$ (d) $0.005\ 3$?
    \vspace*{5cm}
    
    \item[\textbf{21.}] El año tropical, el intervalo desde un equinoccio de primavera hasta el siguiente, es la base para el calendario. Contiene $365.242\ 199$ días. Encuentre el número de segundos en un año tropical.
    \vspace*{5cm}
    
    \item[\textbf{22.}] \textbf{Problema de repaso.} La densidad promedio del planeta Urano es $1.27 \times 10^3 \text{ kg/m}^3$. La proporción de la masa de Neptuno con la de Urano es $1.19$. La proporción del radio de Neptuno con el de Urano es $0.969$. Encuentre la densidad promedio de Neptuno.
    \vspace*{6cm}
    
    \item[\textbf{23.}] \textbf{Problema de repaso.} En un estacionamiento universitario, el número de automóviles ordinarios es mayor que el de vehículos deportivos en $94.7\%$. La diferencia entre ambos números es $18$. Encuentre el número de vehículos deportivos en el estacionamiento.
    \vspace*{5cm}
    
    \item[\textbf{24.}] \textbf{Problema de repaso.} Encuentre todos los ángulos $\theta$ entre $0$ y $360^\circ$ para los cuales la proporción de $\operatorname{sen} \theta$ con $\cos \theta$ sea $-3.00$.
    \vspace*{5cm}
    
    \item[\textbf{25.}] \textbf{Problema de repaso.} La proporción del número de pericos que visita un comedero de aves y el número de aves más interesantes es de $2.25$. Una mañana, cuando $91$ aves visitan el comedero, ¿cuál es el número de pericos?
    \vspace*{5cm}
    
    \item[\textbf{26.}] \textbf{Problema de repaso.} Pruebe que una solución de la ecuación $2.00x^4 - 3.00x^3 + 5.00x = 70.0$ es $x = -2.22$.
    \vspace*{5cm}
    
    \item[\textbf{27.}] \textbf{Problema de repaso.} A partir del conjunto de ecuaciones
    \begin{align*}
        p &= 3q \\
        pr &= qs \\
        \frac{1}{2}pr^2 + \frac{1}{2}qs^2 &= \frac{1}{2}qt^2
    \end{align*}
    que contienen las incógnitas $p, q, r, s$ y $t$, encuentre el valor de la proporción de $t$ con $r$.
    \vspace*{6cm}
    
    \item[\textbf{28.}] \textbf{Problema de repaso.} La figura P1.28 muestra a los estudiantes que analizan la conducción térmica de la energía en bloques cilíndricos de hielo. Este proceso se describe por la ecuación
    $$ \frac{Q}{\Delta t} = \frac{k \pi d^2 (T_h - T_c)}{4L} $$
    Para control experimental, en estos ensayos todas las cantidades, excepto $d$ y $\Delta t$, son constantes. (a) Si $d$ se hace tres veces más grande, ¿la ecuación predice que $\Delta t$ se hará más grande o más pequeña? ¿En qué factor? (b) ¿Qué patrón de proporcionalidad de $\Delta t$ a $d$ predice la ecuación? (c) Para mostrar esta proporcionalidad como una línea recta en una gráfica, ¿qué cantidades debe graficar en los ejes horizontal y vertical? (d) ¿Qué expresión representa la pendiente teórica de esta gráfica?
    \vspace*{10cm}
\end{itemize}

\newpage

\subsection*{PROBLEMAS ADICIONALES}

\begin{itemize}
    \item[\textbf{29.}] En una situación en que los datos se conocen a tres cifras significativas, se escribe $6.379 \text{ m} = 6.38 \text{ m}$ y $6.374 \text{ m} = 6.37 \text{ m}$. Cuando un número termina en $5$, arbitrariamente se elige escribir $6.375 \text{ m} = 6.38 \text{ m}$. Igual se podría escribir $6.375 \text{ m} = 6.37 \text{ m}$, ``redondeando hacia abajo'' en lugar de ``redondear hacia arriba'', porque el número $6.375$ se cambiaría por incrementos iguales en ambos casos. Ahora considere una estimación del orden de magnitud en la cual los factores de cambio, más que los incrementos, son importantes. Se escribe $500 \text{ m} \sim 10^3 \text{ m}$ porque $500$ difiere de $100$ por un factor de $5$, mientras difiere de $1000$ solo por un factor de $2$. Escriba $437 \text{ m} \sim 10^3 \text{ m}$ y $305 \text{ m} \sim 10^2 \text{ m}$. ¿Qué distancia difiere de $100 \text{ m}$ y de $1000 \text{ m}$ por factores iguales de modo que lo mismo se podría escoger representar su orden de magnitud como $\sim 10^2 \text{ m}$ o como $\sim 10^3 \text{ m}$?
    \vspace*{7cm}
    
    \item[\textbf{30.}] (a) ¿Cuál es el orden de magnitud del número de microorganismos en el tracto intestinal humano? Una escala de longitud bacteriana típica es de $10^{-6} \text{ m}$. Estime el volumen intestinal y suponga que $1\%$ del mismo está ocupada por una bacteria. (b) ¿El número de bacterias indica si éstas son beneficiosas, peligrosas o neutras para el cuerpo humano? ¿Qué funciones podrían tener?
    \vspace*{8cm}
    
    \item[\textbf{31.}] La distancia del Sol a la estrella más cercana es casi de $4 \times 10^{16} \text{ m}$. La Vía Láctea (figura P1.31) es, en términos aproximados, un disco de $\sim 10^{21} \text{ m}$ de diámetro y $\sim 10^{19} \text{ m}$ de grosor. Encuentre el orden de magnitud del número de estrellas en la Vía Láctea. Considere representativa la distancia entre el Sol y el vecino más cercano.
    \vspace*{8cm}
    
    \item[\textbf{32.}] ¿Por qué no es posible la siguiente situación? En un esfuerzo para aumentar el interés en un programa de televisión, a cada ganador semanal se le ofrece un premio de bonificación de $\$1$ millón adicional si él o ella pueden contar personalmente la cantidad exacta de un paquete de billetes de un dólar. El ganador deberá realizar esta tarea bajo la supervisión de los ejecutivos del programa de televisión y en una semana laboral de $40$ horas. Para consternación de los productores del programa, la mayoría de los participantes tienen éxito en el desafío.
    \vspace*{7cm}
    
    \item[\textbf{33.}] Las bacterias y otros procariotas se encuentran bajo tierra, en el agua y en el aire. Una micra ($10^{-6} \text{ m}$) es una escala de longitud característica asociada con estos microbios. (a) Estime el número total de bacterias y otros procariotas sobre la Tierra. (b) Estime la masa total de todos estos microbios.
    \vspace*{8cm}
    
    \item[\textbf{34.}] Un cascarón esférico tiene un radio externo de $2.60 \text{ cm}$ y uno interno de $a$. La pared del cascarón tiene grosor uniforme y está hecho de un material cuya densidad es de $4.70 \text{ g/cm}^3$. El espacio interior del cascarón está lleno con un líquido que tiene una densidad de $1.23 \text{ g/cm}^3$. (a) Encuentre la masa $m$ de la esfera, incluidos sus contenidos, como función de $a$. (b) ¿Para qué valor de $a$ tiene $m$ su máximo valor posible? (c) ¿Cuál es esta masa máxima? (d) Explique si el valor de la parte (c) concuerda con el resultado de un cálculo directo de la masa de una esfera sólida de densidad uniforme hecha del mismo material que el cascarón? (e) ¿Qué pasaría si? En el inciso (a), ¿la respuesta cambiaría si la pared interior del cascarón no fuese concéntrica con la pared exterior?
    \vspace*{12cm}
    
    \item[\textbf{35.}] Se sopla aire hacia adentro de un globo esférico de modo que, cuando su radio es de $6.50 \text{ cm}$, este aumenta con una rapidez de $0.900 \text{ cm/s}$. (a) Encuentre la rapidez a la que aumenta el volumen del globo. (b) Si dicha rapidez de flujo volumétrico de aire que entra al globo es constante, ¿en qué proporción aumentará el radio cuando este es de $13.0 \text{ cm}$? (c) Explique físicamente por qué la respuesta del inciso (b) es mayor o menor que $0.9 \text{ cm/s}$, si es diferente.
    \vspace*{8cm}
    
    \item[\textbf{36.}] En física es importante usar aproximaciones matemáticas. (a) Demuestre que, para ángulos pequeños ($< 20^\circ$),
    $$ \tan \alpha \approx \operatorname{sen} \alpha \approx \alpha = \frac{\pi \alpha'}{180^\circ} $$
    donde $\alpha$ está en radianes y $\alpha'$ en grados. (b) Use una calculadora para encontrar el ángulo más grande para el que tan $\alpha$ se pueda aproximar a $\alpha$ con un error menor de $10.0$ por ciento.
    \vspace*{8cm}
    
    \item[\textbf{37.}] El consumo de gas natural por una compañía satisface la ecuación empírica $V = 1.50t + 0.008\ 00t^2$, donde $V$ es el volumen en millones de pies cúbicos y $t$ es el tiempo en meses. Exprese esta ecuación en unidades de pies cúbicos y segundos. Suponga un mes de $30.0$ días.
    \vspace*{8cm}
    
    \item[\textbf{38.}] Una mujer que desea conocer la altura de una montaña mide el ángulo de elevación de la cima, que es de $12.0^\circ$. Después de caminar $1.00 \text{ km}$ más cerca de la montaña a nivel del suelo, encuentra que el ángulo es $14.0^\circ$. (a) Realice un dibujo del problema, ignorando la altura de los ojos de la mujer por encima del suelo. \textit{Sugerencia:} Utilice dos triángulos. (b) Usando el símbolo $y$ para representar la altura de la montaña y el símbolo $x$ para representar la distancia original de la mujer a la montaña, etiquete la imagen. (c) Utilizando la imagen etiquetada, escriba dos ecuaciones trigonométricas que relacionen las dos variables seleccionadas. (d) Determine la altura $y$.
    \vspace*{10cm}
\end{itemize}

\newpage

\subsection*{PROBLEMA DE DESAFÍO}

\begin{itemize}
    \item[\textbf{39.}] Una mujer se encuentra a una distancia horizontal $x$ de una montaña y mide el ángulo $\theta$ de elevación de la cima sobre la horizontal. Después de caminar acercándose a la montaña una distancia $d$ a nivel del suelo, encuentra que el ángulo es $\phi$. Encuentre una ecuación general para la altura $y$ de la montaña en términos de $d$, $\phi$ y $\theta$; ignore la altura de los ojos al suelo.
    \vspace*{10cm}
\end{itemize}

\end{document}
